\scanpage{643}
\mspage{31}

\uline{Scholie 2.10} Soit $(W,C,F)$ un triple de Quillen
clos dans $M$. Pour \deleted{\unsure{un}} $I$ dans Cat, pour que
$(W(I),C\langle I\rangle,F(I))$ soit un triple de Quillen
clos pour $\unsure{M(I)}$, il f et s. que les
couples $C\langle I\rangle\longleftrightarrow TF(I)$ et $TC\langle I\rangle\longleftrightarrow F(I)$
soient factorisants. Supposons que cette
condition soit satisfaite pour $I$ et pour $I'$,
et soit $u\colon I\longrightarrow I'$ une flèche \deleted{\unsure{dans ?}}\inserted{de}
$I$ dans $I'$. Alors les assertions de 2.3. $2^\circ$
sont valables, en particulier on trouve un
foncteur $u_!^W\colon\operatorname{Ho}_W(I)\unsure{-|\longrightarrow}\operatorname{Ho}_W(I')$, adjoint
à gauche de $u^*_W$.

% slanted
\margin{Dire aussi que le calcul des fibres de $u_!^W(X)$, prévu dans Der 5 bis, est OK.}

Ces résultats se
dualisent aussitôt (par $M\unsure{\mapsto|-} M^{\unsure{o|\circ}}, I\unsure{-|\mapsto}I^{o}$)
pour les triples $(W(I),C(I),F[I])$,
et pour la construction de $u_*^W$, adjoint
à droite de $u^*_W$. En particulier, pour que
le triple précédent soit un triple de Quillen
(forcément clos) il f. et s. que les deux
couples $C(I)\longleftrightarrow TF[I]$ et $TC(I)\longleftrightarrow F[I]$
soient factorisants ; et quand cette condition
est satisfaite pour $I$ et pour $I'$, \deleted{\unsure{et}}(pour une
flèche $u\colon I\mathrel{-} I'$ dans Cat donnée), alors
$u_*^W$ existe et se construit dualement de
la construction de $u_!^W$ explicitée dans
2.3. $2^\circ$.
% archive sheet number: 642
